% W4i39_ThirdPeak_Resolution.tex
% DFD 17-Agent Research Programme — Iteration 39 (i39)
% Agents 6 (Vector Extension), 7 (SymBoltz.jl), 8 (Nonlinear Mode Coupling),
%         9 (Baryon Feedback), 10 (Recombination)
% Iteration 8/20 of Quantum Gravity Cycle (i32-i51)
% PRIME DIRECTIVE: Solve the CMB Third Peak Crisis
% Date: 2026-03-22

\documentclass[11pt,a4paper]{article}
\usepackage[utf8]{inputenc}
\usepackage{amsmath,amssymb,amsfonts,amsthm}
\usepackage{hyperref}
\usepackage{booktabs}
\usepackage{longtable}
\usepackage{geometry}
\usepackage{xcolor}
\geometry{margin=2.5cm}

\newtheorem{theorem}{Theorem}
\newtheorem{proposition}{Proposition}
\newtheorem{lemma}{Lemma}
\newtheorem{corollary}{Corollary}
\newtheorem{definition}{Definition}
\newtheorem{remark}{Remark}

\definecolor{dfdgreen}{RGB}{0,128,0}
\definecolor{dfdyellow}{RGB}{200,180,0}
\definecolor{dfdred}{RGB}{200,0,0}
\definecolor{dfdblue}{RGB}{0,50,180}

\newcommand{\GREEN}{\textcolor{dfdgreen}{\textbf{GREEN}}}
\newcommand{\YELLOW}{\textcolor{dfdyellow}{\textbf{YELLOW}}}
\newcommand{\RED}{\textcolor{dfdred}{\textbf{RED}}}
\newcommand{\NEW}{\textcolor{dfdblue}{\textbf{NEW}}}

\title{DFD i39: CMB Third Peak Crisis --- Resolution Paths\\[6pt]
\large Agents 6 (Vector), 7 (SymBoltz), 8 (Nonlinear), 9 (Feedback), 10 (Recomb.)\\[6pt]
\normalsize Iteration 8/20 of Quantum Gravity Cycle (i32--i51)\\
PRIME DIRECTIVE: Solve the CMB Third Peak Crisis}
\author{DFD 17-Agent Research Programme}
\date{2026-03-22}

\begin{document}
\maketitle
\tableofcontents
\newpage

%=============================================================================
\section{Agent 6: Vector Extension --- Minimal Cost Assessment}
\label{sec:agent6}
%=============================================================================

\subsection{Mandate}

If DFD \emph{must} add a vector field, define the minimal extension. Assess what it costs (new parameters, broken predictions, theoretical complexity).

\subsection{Key Papers}

\begin{itemize}
\item Skordis \& Z{\l}o\'{s}nik (2021), PRL 127, 161302 --- AeST theory
\item Jacobson \& Mattingly (2001), PRD 64, 024028 --- Einstein-aether theory
\item Durakovi\'{c} et al.\ (2025), arXiv:2503.11174 --- mimetic gravity connection
\item Zlosnik et al.\ (2007), PRD 75, 044017 --- vector-tensor MOND
\item Sanders (2005), MNRAS 363, 459 --- TeVeS cosmology
\item Blanchet et al.\ (2024), JCAP 2024, 040 --- relativistic khronon theory
\end{itemize}

\subsection{The Minimal Vector Extension of DFD}

\subsubsection{What Is Needed}

From Agent 1's analysis, the third peak requires a gravitational degree of freedom that:
\begin{enumerate}
\item Does not couple directly to the baryon stress-energy (so it doesn't oscillate)
\item Grows or remains constant during radiation domination
\item Provides gravitational potential wells at scales $k \sim 0.01$--$0.1$ Mpc$^{-1}$
\item Transitions to MOND-like behavior on galaxy scales
\end{enumerate}

A unit-timelike vector field $A^\mu$ with $g_{\mu\nu}A^\mu A^\nu = -1$ satisfies all four requirements in AeST.

\subsubsection{The DFD + Vector Action}

The minimal extension adds:
\begin{equation}
S_{\rm DFD+V} = S_{\rm DFD} + \frac{1}{16\pi G}\int d^4x\,\sqrt{-g}\,\mathcal{L}_A
\end{equation}

Following the AeST structure:
\begin{equation}
\mathcal{L}_A = \frac{K_B}{2}F_{\mu\nu}F^{\mu\nu} + \frac{K_S}{2}\left(\nabla_\mu A^\mu\right)^2 + \lambda(g_{\mu\nu}A^\mu A^\nu + 1) + \mathcal{L}_{\rm int}(\psi, A^\mu)
\end{equation}

where:
\begin{itemize}
\item $F_{\mu\nu} = \nabla_\mu A_\nu - \nabla_\nu A_\mu$ (Maxwell-like kinetic term)
\item $K_B, K_S$ are dimensionless coupling constants
\item $\lambda$ is a Lagrange multiplier enforcing the unit-timelike constraint
\item $\mathcal{L}_{\rm int}$ couples the vector to DFD's scalar $\psi$
\end{itemize}

\subsubsection{New Parameters}

\begin{center}
\begin{tabular}{lll}
\toprule
\textbf{Parameter} & \textbf{Role} & \textbf{Constraint} \\
\midrule
$K_B$ & Vector kinetic mixing & CMB peak heights \\
$K_S$ & Vector expansion coupling & Growth rate \\
$\alpha_{A\psi}$ & Vector-scalar coupling & MOND-to-CMB transition \\
\bottomrule
\end{tabular}
\end{center}

\textbf{Three new parameters.} This is the minimum. AeST has $\sim$5--7 free parameters for the cosmological fit.

\subsubsection{The Interaction Term}

Durakovi\'{c} et al.\ (2025) showed that any relativistic MOND theory with a unit-timelike vector can be embedded in a conformal/disformal-invariant framework. The natural interaction is:
\begin{equation}
\mathcal{L}_{\rm int} = f(\psi)\left(A^\mu\nabla_\mu\psi\right)^2
\end{equation}

This couples the vector to the scalar through the projection of $\nabla\psi$ along the aether direction. In the MOND limit, this reduces to:
\begin{equation}
\mathcal{L}_{\rm int} \to f(\psi)\dot{\psi}^2
\end{equation}
which modifies the scalar kinetic term in a direction-dependent way.

\subsection{Impact on Existing DFD Predictions}

\subsubsection{What Breaks}

\begin{enumerate}
\item \textbf{$c_{\rm GW} = c$:} The vector kinetic term can modify the gravitational wave speed. Need $K_B + K_S = 0$ (or fine-tuned) to maintain $c_{\rm GW} = c$ exactly. AeST achieves this.

\item \textbf{Black hole entropy:} The vector field modifies the horizon structure. The Wald entropy calculation (i35--i37) would need revision. Estimated correction: $\delta S_{\rm BH}/S_{\rm BH} \sim K_B$.

\item \textbf{Born rule / einselection (i34):} The constraint algebra changes with a vector field. The Born rule derivation from the $\psi$ constraint field may be modified.

\item \textbf{BMV prediction:} The tabletop gravity experiment prediction depends on $G_{\rm eff}$ in the Newtonian regime, where the vector field should decouple. Estimated impact: negligible.

\item \textbf{Lab predictions (P1--P15):} All depend on the scalar sector. The vector field decouples at lab scales ($g \gg a_0$). Impact: negligible.
\end{enumerate}

\subsubsection{What Survives}

\begin{itemize}
\item All 15 patent predictions (lab-scale, $g \gg a_0$, vector decouples)
\item Galaxy rotation curves (MOND regime preserved by construction)
\item Tully-Fisher relation (MOND)
\item GW speed $= c$ (with AeST-like constraint)
\item $D_0 = 0.017$ (scalar sector parameter)
\item Bare $\Lambda$ (expansion history unchanged)
\end{itemize}

\subsection{Theoretical Cost Assessment}

\begin{center}
\begin{tabular}{lcl}
\toprule
\textbf{Cost} & \textbf{Severity} & \textbf{Notes} \\
\midrule
New field & HIGH & Adds 3 propagating DOF \\
New parameters ($K_B, K_S, \alpha_{A\psi}$) & MEDIUM & 3 new dimensionless constants \\
$c_{\rm GW}$ constraint & LOW & $K_B + K_S = 0$ (follows AeST) \\
BH entropy revision & MEDIUM & Wald entropy needs recalculation \\
Born rule revision & HIGH & Core QG result at risk \\
Aesthetic/minimality & HIGH & DFD's claim of ``just one scalar'' dies \\
\bottomrule
\end{tabular}
\end{center}

\subsection{Agent 6 Conclusion}

\begin{center}
\fbox{\parbox{0.9\textwidth}{
\textbf{Finding 6.1:} The minimal vector extension adds 3 new parameters ($K_B, K_S, \alpha_{A\psi}$) and one new field ($A^\mu$, unit-timelike vector with 3 DOF).\\[6pt]
\textbf{Finding 6.2:} Lab-scale predictions (all 15 patents, BMV, GW echoes) are unaffected because the vector decouples at $g \gg a_0$.\\[6pt]
\textbf{Finding 6.3:} Quantum gravity results (Born rule, einselection, BH entropy) may require revision. This is the highest theoretical cost.\\[6pt]
\textbf{Finding 6.4:} The vector extension converts DFD into an AeST-like theory. It is viable for the CMB but sacrifices DFD's minimality.\\[6pt]
\textbf{Finding 6.5:} Durakovi\'{c} et al.\ (2025) show the vector can be embedded in a conformal/disformal-invariant framework, suggesting it may be a natural extension rather than an ad hoc addition.\\[6pt]
\textbf{STATUS:} \YELLOW{} --- Vector extension works for the CMB but has significant theoretical costs. The Born rule derivation must be re-examined.
}}
\end{center}

\newpage
%=============================================================================
\section{Agent 7: SymBoltz.jl --- Numerical CMB Computation for DFD}
\label{sec:agent7}
%=============================================================================

\subsection{Mandate}

Define exactly what SymBoltz.jl would need to compute DFD's CMB. Could the third-peak deficit be an artifact of analytic approximation?

\subsection{Key Papers}

\begin{itemize}
\item Hersle (2025), A\&A, arXiv:2509.24740 --- SymBoltz.jl
\item Zack Li, Bolt.jl --- differentiable Boltzmann code
\item Lewis et al.\ (2000), ApJ 538, 473 --- CAMB
\item Lesgourgues (2011), arXiv:1104.2932 --- CLASS
\item Zumalac\'{a}rregui et al.\ (2017), JCAP 1708, 019 --- hi\_class (modified gravity)
\item Skordis et al.\ (2006), PRL 96, 011301 --- TeVeS CMB computation
\end{itemize}

\subsection{SymBoltz.jl Capabilities}

SymBoltz.jl (Hersle 2025) is a symbolic-numeric Einstein-Boltzmann solver in Julia with key features:

\begin{enumerate}
\item \textbf{Symbolic equation specification:} Models defined in high-level symbolic format, compiled to fast numerical functions.
\item \textbf{Approximation-free:} No tight-coupling, ultrarelativistic fluid, or radiation streaming approximations.
\item \textbf{Differentiable:} Exact derivatives via automatic differentiation.
\item \textbf{Extensible:} Modified models can change constraint equations into differential equations.
\end{enumerate}

Version 1.0.0 includes: flat Newtonian gauge, GR, CDM, photons, baryons, RECFAST recombination, massless and massive neutrinos, $\Lambda$, $w_0 w_a$ dark energy, distances, matter and CMB spectra.

\subsection{DFD Implementation Requirements}

\subsubsection{Modified Poisson Equation}

In standard $\Lambda$CDM:
\begin{equation}
k^2\Phi = -4\pi G a^2 \sum_i \rho_i\delta_i
\end{equation}

In DFD, this becomes:
\begin{equation}
k^2\Phi = -4\pi G_{\rm eff}(k,\eta) a^2 \rho_b\delta_b - 4\pi G a^2(\rho_\gamma\delta_\gamma + \rho_\nu\delta_\nu) + S_\psi(k,\eta)
\end{equation}

where:
\begin{equation}
G_{\rm eff}(k,\eta) = G_N\cdot\nu\left(\frac{k^2|\Phi|}{4\pi G_N a^2\rho_b\delta_b \cdot a_0}\right)
\end{equation}

and $S_\psi$ is the scalar field source term:
\begin{equation}
S_\psi = -\frac{D_0}{M_P}\left(\ddot{\delta\psi} + 2\mathcal{H}\dot{\delta\psi} + k^2\delta\psi + V''(\psi)\delta\psi\right)
\end{equation}

\textbf{CRITICAL:} This is a \emph{nonlinear} equation ($G_{\rm eff}$ depends on $\Phi$, which depends on $G_{\rm eff}$). Standard Boltzmann solvers linearize. SymBoltz.jl can handle this because it solves the full ODE system without approximation switching.

\subsubsection{Scalar Field Evolution}

The $\psi$ perturbation equation:
\begin{equation}
\ddot{\delta\psi} + 2\mathcal{H}\dot{\delta\psi} + \left(k^2 + a^2 V''(\bar{\psi})\right)\delta\psi = \frac{D_0}{M_P}a^2\delta T_b + \text{metric source terms}
\end{equation}

where $\delta T_b$ includes baryon density and pressure perturbations.

\subsubsection{Disformal Metric Perturbations}

The perturbed disformal metric:
\begin{equation}
\delta\tilde{g}_{00} = -a^2\left[2e^{2\bar{\psi}}\Psi + 2e^{2\bar{\psi}}\delta\psi + D'(\bar{\chi})\delta\chi\dot{\bar{\psi}}^2 + 2D(\bar{\chi})\dot{\bar{\psi}}\delta\dot{\psi}\right]
\end{equation}
\begin{equation}
\delta\tilde{g}_{0i} = a^2\left[e^{2\bar{\psi}}\partial_i B + D(\bar{\chi})\dot{\bar{\psi}}\partial_i\delta\psi\right]
\end{equation}
\begin{equation}
\delta\tilde{g}_{ij} = a^2 e^{2\bar{\psi}}\left[(1-2\Phi)\delta_{ij} + 2\delta\psi\delta_{ij}\right]
\end{equation}

\subsubsection{What SymBoltz.jl Needs}

\begin{enumerate}
\item The background equations for $\bar{\psi}(\eta)$
\item The perturbation equations for $\delta\psi(k,\eta)$
\item The modified Poisson equation with $G_{\rm eff}(k,\eta)$
\item The modified geodesic equations for photons on $\tilde{g}_{\mu\nu}$
\item Initial conditions for $\delta\psi$ (adiabatic or isocurvature)
\item The nonlinear solver for the self-consistent $G_{\rm eff} \leftrightarrow \Phi$ loop
\end{enumerate}

\subsection{Could the Deficit Be an Artifact?}

\subsubsection{Sources of Analytic Error}

Agent 3's calculation used:
\begin{itemize}
\item Scale-by-scale MOND enhancement (ignoring mode coupling)
\item Instantaneous $G_{\rm eff}$ (ignoring time evolution through recombination)
\item Baryon-only potential (ignoring $\delta\psi$ backreaction)
\item Static MOND interpolation function (ignoring time-dependent $a_0$)
\end{itemize}

\subsubsection{Potential for the Deficit to Be Smaller}

\textbf{(a) Time-dependent $G_{\rm eff}$:} During recombination ($z \sim 1100 \to 900$), the baryon perturbation evolves, so $g_N$ changes, so $G_{\rm eff}$ changes. If $G_{\rm eff}$ is larger at early times (deeper MOND) and smaller at late times (approach to Newtonian), the time-averaged enhancement could be more uniform across scales.

\textbf{Estimate:} The recombination width is $\Delta z \sim 200$. Over this time, $g_N$ changes by $\sim 20\%$ at the third-peak scale. The effect on $G_{\rm eff}$: $\delta\nu/\nu \sim 10\%$. This does not close the factor-of-5 gap.

\textbf{(b) $\delta\psi$ backreaction:} The scalar perturbation sources additional potential. But Agent 1 showed $\delta\psi$ oscillates with baryons, so it does not help the third peak preferentially.

\textbf{(c) Mode coupling:} See Agent 8.

\textbf{HONEST ASSESSMENT:} The analytic estimate may be off by factors of $\sim 2$, but the direction (third peak suppressed relative to first) is robust. A full numerical computation is unlikely to change a factor-of-5 deficit into agreement.

\subsection{Agent 7 Conclusion}

\begin{center}
\fbox{\parbox{0.9\textwidth}{
\textbf{Finding 7.1:} SymBoltz.jl is the ideal tool for DFD's CMB computation. Its symbolic-numeric, approximation-free architecture can handle DFD's nonlinear $G_{\rm eff}$.\\[6pt]
\textbf{Finding 7.2:} Implementation requires: modified Poisson equation, scalar field evolution, disformal metric perturbations, and a nonlinear solver for the $G_{\rm eff} \leftrightarrow \Phi$ loop.\\[6pt]
\textbf{Finding 7.3:} The third-peak deficit is unlikely to be an analytic artifact. Time-dependent corrections are $\sim 10\%$, insufficient to close the factor-of-5 gap identified by Agent 3.\\[6pt]
\textbf{Finding 7.4:} A full SymBoltz.jl implementation should be the next concrete step, as it would definitively quantify the deficit and determine whether \emph{any} DFD parameter choice can fit the CMB.\\[6pt]
\textbf{STATUS:} \YELLOW{} --- SymBoltz.jl computation is essential but unlikely to save DFD's current scalar-only CMB. It would definitively quantify the problem.
}}
\end{center}

\newpage
%=============================================================================
\section{Agent 8: Nonlinear Mode Coupling}
\label{sec:agent8}
%=============================================================================

\subsection{Mandate}

Assess whether nonlinear mode coupling in DFD's MOND regime could increase the third peak beyond linear-theory predictions.

\subsection{Key Papers}

\begin{itemize}
\item Bartolo et al.\ (2004), Phys.\ Rep.\ 402, 103 --- non-Gaussianity from inflation
\item Pitrou et al.\ (2010), JCAP 1007, 003 --- second-order CMB
\item Hu \& Okamoto (2002), ApJ 574, 566 --- CMB lensing
\item Pettinari et al.\ (2013), JCAP 1304, 003 --- SONG (second-order code)
\item Beneke et al.\ (2011), JCAP 1110, 016 --- second-order perturbations
\item Mollerach \& Matarrese (1997), PRD 56, 4494 --- CMB beyond linear
\end{itemize}

\subsection{The Case for Nonlinear Enhancement}

\subsubsection{Linear Theory Breakdown in MOND}

In MOND, the effective gravitational constant is $G_{\rm eff} = G_N \nu(g_N/a_0)$, which is a \emph{nonlinear} function of the density perturbation. Expanding:
\begin{equation}
G_{\rm eff}(\delta) = G_{\rm eff}(\bar{\delta}) + G_{\rm eff}'(\bar{\delta})\delta + \frac{1}{2}G_{\rm eff}''(\bar{\delta})\delta^2 + \ldots
\end{equation}

The second-order term generates mode coupling:
\begin{equation}
\Phi^{(2)}(\mathbf{k}) = \int \frac{d^3k'}{(2\pi)^3}\,\mathcal{K}(\mathbf{k}',\mathbf{k}-\mathbf{k}')\,\delta(\mathbf{k}')\delta(\mathbf{k}-\mathbf{k}')
\end{equation}

where $\mathcal{K}$ is the MOND mode-coupling kernel, which is larger than the GR kernel because $G_{\rm eff}''$ is larger.

\subsubsection{Can Mode Coupling Boost the Third Peak?}

The third peak ($\ell \sim 800$, $k \sim 0.06$ Mpc$^{-1}$) can receive power from two coupled modes:
\begin{equation}
k_3 = k_1 + k_2
\end{equation}

For example, the first peak mode ($k_1 \sim 0.02$) coupling with a mode near the second peak ($k_2 \sim 0.04$). The second-order contribution:
\begin{equation}
\Phi^{(2)}(k_3) \sim \mathcal{K}(k_1, k_2)\,\delta(k_1)\delta(k_2) \sim \mathcal{K}\,\times\,10^{-5}\times10^{-5} = \mathcal{K}\times10^{-10}
\end{equation}

The first-order contribution at the third peak: $\Phi^{(1)}(k_3) \sim G_{\rm eff}(k_3)\times\delta_b(k_3) \sim 2G_N \times 10^{-5} \sim 10^{-5}$.

\textbf{Ratio:}
\begin{equation}
\frac{\Phi^{(2)}}{\Phi^{(1)}} \sim \frac{\mathcal{K}\times10^{-10}}{10^{-5}} = \mathcal{K}\times10^{-5}
\end{equation}

Even if $\mathcal{K} \sim 10^3$ (much larger than GR), the second-order correction is $\sim 10^{-2}$, i.e., 1\%. This is completely negligible compared to the factor-of-5 deficit.

\subsubsection{Higher-Order Corrections}

The $n$-th order correction scales as $(\delta_b)^{n-1} \sim (10^{-5})^{n-1}$. Even with MOND-enhanced kernels, no perturbative order can bridge the factor-of-5 gap.

\subsection{Non-Perturbative Effects}

Could there be non-perturbative MOND effects that evade the perturbative argument? This is possible only if the MOND interpolation function has singularities or rapid transitions near $y \sim 1$. For DFD's $\nu(y) = (1+\sqrt{1+4/y})/2$, the function is smooth everywhere. There are no non-perturbative enhancements.

\subsection{Agent 8 Conclusion}

\begin{center}
\fbox{\parbox{0.9\textwidth}{
\textbf{Finding 8.1:} Nonlinear mode coupling in DFD's MOND regime generates second-order corrections $\sim 10^{-2}$ relative to first order. This is completely negligible.\\[6pt]
\textbf{Finding 8.2:} The perturbative expansion converges rapidly ($\delta_b \sim 10^{-5}$ at recombination). No perturbative order can bridge the factor-of-5 gap.\\[6pt]
\textbf{Finding 8.3:} DFD's MOND interpolation function is smooth, ruling out non-perturbative enhancements.\\[6pt]
\textbf{STATUS:} \RED{} --- Nonlinear mode coupling cannot help. The deficit is a linear-theory problem.
}}
\end{center}

\newpage
%=============================================================================
\section{Agent 9: Baryon Feedback --- The $\psi$-Baryon Loop}
\label{sec:agent9}
%=============================================================================

\subsection{Mandate}

Assess whether the baryon-$\psi$ feedback loop at recombination is stronger than estimated, potentially boosting the third peak.

\subsection{Key Papers}

\begin{itemize}
\item Hu \& Sugiyama (1996), ApJ 471, 542 --- photon-baryon driving
\item Sakstein et al.\ (2025), arXiv:2501.15298 --- scalar-tensor gravity and BAO
\item Bruni et al.\ (2024), arXiv:2407.15832 --- nonminimally coupled gravity
\item Amendola (2000), PRD 62, 043511 --- coupled quintessence
\item Pettorino \& Baccigalupi (2008), PRD 77, 103003 --- coupled dark energy
\item Dalang \& Lombriser (2025), arXiv:2602.19576 --- K-essence Boltzmann dynamics
\end{itemize}

\subsection{The Baryon-$\psi$ Feedback Loop}

\subsubsection{The Coupling Chain}

In DFD:
\begin{enumerate}
\item Baryons create potential $\Phi \to$ sources $\delta\psi$ via $\Box\psi = (D_0/M_P)T_b$
\item $\delta\psi$ modifies the effective metric $\tilde{g}_{\mu\nu} \to$ baryons feel different gravitational force
\item Modified baryon dynamics $\to$ changed $\delta_b \to$ changed $\Phi$
\item Repeat
\end{enumerate}

The feedback amplification factor:
\begin{equation}
\mathcal{A} = \frac{1}{1 - D_0\cdot\mathcal{G}(k,\eta)}
\end{equation}

where $\mathcal{G}(k,\eta)$ is the Green's function for the $\psi$-response to a baryon source. If $D_0\mathcal{G} \to 1$, the amplification diverges (resonance).

\subsubsection{Estimating $\mathcal{G}$}

The Green's function for $\psi$ on the cosmological background:
\begin{equation}
\mathcal{G}(k,\eta) \sim \frac{1}{k^2 + m_\psi^2 a^2}\cdot\frac{a^2\rho_b}{M_P}
\end{equation}

At the third-peak scale ($k \sim 0.06$ Mpc$^{-1}$), with $m_\psi \sim H_0 \sim 10^{-33}$ eV (negligible at recombination):
\begin{equation}
\mathcal{G} \sim \frac{a^2\rho_b}{k^2 M_P} \sim \frac{(10^{-3})^2 \times 6.5\times10^{-22}\,\text{g/cm}^3}{(0.06\,\text{Mpc}^{-1})^2 \times M_P}
\end{equation}

Converting units ($\rho_b$ in natural units $\sim 10^{-90}\,M_P^4$, $k^2 \sim 10^{-58}\,M_P^2$):
\begin{equation}
\mathcal{G} \sim \frac{10^{-6}\times10^{-90}}{10^{-58}} M_P = 10^{-38} M_P
\end{equation}

So $D_0\mathcal{G} \sim 0.017 \times 10^{-38} M_P / M_P = 10^{-40}$.

\textbf{The feedback amplification is negligible:}
\begin{equation}
\mathcal{A} = \frac{1}{1 - 10^{-40}} \approx 1 + 10^{-40}
\end{equation}

\subsubsection{Why So Small?}

The coupling constant $D_0/M_P \sim 10^{-2}/M_P$ is extremely weak in Planck units. The baryon-$\psi$ coupling is gravitational-strength, and at recombination the baryon density is far below the Planck density. The feedback loop is completely negligible.

\subsection{Could a Stronger Coupling Help?}

If we allowed $D_0 \sim 1$ (strong coupling), the amplification would be $\sim 10^{-38}$ --- still negligible. Even $D_0 \sim M_P$ (which would destroy all DFD predictions) gives $\mathcal{A} \sim 1 + 10^{-38}$.

The reason: the feedback loop involves two gravitational-strength couplings (baryon$\to\psi$ and $\psi\to$metric), each suppressed by $M_P^{-1}$. The product is $M_P^{-2}$, which at cosmological densities is infinitesimal.

\subsection{Agent 9 Conclusion}

\begin{center}
\fbox{\parbox{0.9\textwidth}{
\textbf{Finding 9.1:} The baryon-$\psi$ feedback amplification at recombination is $\mathcal{A} \approx 1 + 10^{-40}$. Completely negligible.\\[6pt]
\textbf{Finding 9.2:} The feedback is weak because it involves two $M_P^{-1}$-suppressed couplings multiplied by cosmological-density matter.\\[6pt]
\textbf{Finding 9.3:} No reasonable modification of $D_0$ can make the feedback significant without destroying DFD's other predictions.\\[6pt]
\textbf{STATUS:} \RED{} --- Baryon-$\psi$ feedback cannot help. The coupling is gravitationally weak.
}}
\end{center}

\newpage
%=============================================================================
\section{Agent 10: Recombination --- $\psi$-Modified Atomic Physics}
\label{sec:agent10}
%=============================================================================

\subsection{Mandate}

Determine whether DFD's $\psi$-dependent modifications to atomic physics (fine structure constant variation) can shift the third peak.

\subsection{Key Papers}

\begin{itemize}
\item Hart \& Chluba (2020), MNRAS 493, 3255 --- updated RECFAST
\item Planck Collaboration (2016), A\&A 596, A107 --- varying $\alpha$ constraints
\item Sekiguchi \& Takahashi (2021), PRD 103, 083507 --- varying constants and CMB
\item Scoccola et al.\ (2009), PRD 79, 083517 --- varying $\alpha$ recombination
\item L\'{o}pez-Honorez et al.\ (2024), arXiv:2403.02256 --- incompatibility of $\alpha$ variations
\item Murphy et al.\ (2025), arXiv:2512.14441 --- sub-ppm upper limit on $\alpha$ variation
\end{itemize}

\subsection{$\psi$-Dependent Fine Structure Constant}

\subsubsection{The Mechanism}

In DFD, electromagnetic coupling occurs through the effective metric $\tilde{g}_{\mu\nu}$. If $\psi$ couples to the EM field:
\begin{equation}
\mathcal{L}_{\rm EM} = -\frac{1}{4}e^{-2\psi}F_{\mu\nu}F^{\mu\nu}
\end{equation}

This gives an effective fine structure constant:
\begin{equation}
\alpha_{\rm eff} = \alpha_0\,e^{2\psi}
\end{equation}

At recombination ($\bar{\psi} \sim 10^{-7}$):
\begin{equation}
\frac{\Delta\alpha}{\alpha} = 2\bar{\psi}_{\rm rec} \sim 2\times10^{-7}
\end{equation}

\subsubsection{Effect on Recombination}

A change in $\alpha$ modifies:
\begin{enumerate}
\item The hydrogen binding energy: $E_{\rm ion} = 13.6\,\text{eV}\times(\alpha/\alpha_0)^2$
\item The Thomson cross-section: $\sigma_T \propto \alpha^2$
\item The recombination rate: $\alpha_{\rm rec} \propto \alpha^5$
\end{enumerate}

Scoccola et al.\ (2009) showed that $\Delta\alpha/\alpha = +10^{-2}$ shifts the recombination redshift by $\Delta z_{\rm rec} \approx +30$ and modifies peak heights by $\sim 5\%$.

For DFD's $\Delta\alpha/\alpha \sim 2\times10^{-7}$:
\begin{equation}
\Delta z_{\rm rec} \approx 30 \times \frac{2\times10^{-7}}{10^{-2}} = 6\times10^{-4}
\end{equation}
\begin{equation}
\frac{\Delta C_\ell}{C_\ell} \approx 5\% \times \frac{2\times10^{-7}}{10^{-2}} = 10^{-6}
\end{equation}

\textbf{Completely negligible.} The $\alpha$ variation in DFD is far too small to affect the CMB.

\subsubsection{Current Observational Constraints}

L\'{o}pez-Honorez et al.\ (2024) showed that varying $\alpha$ models face a fundamental incompatibility: values needed to affect recombination ($\Delta\alpha/\alpha \gtrsim 10^{-3}$) are in tension with local measurements ($\Delta\alpha/\alpha < 10^{-6}$, Murphy et al.\ 2025).

DFD predicts $\Delta\alpha/\alpha \sim 10^{-7}$, which is safely below all constraints but also far too small to affect the CMB.

\subsection{$\psi$-Dependent Electron Mass}

If $\psi$ also modifies the electron mass:
\begin{equation}
m_e^{\rm eff} = m_e\,e^{\psi}
\end{equation}

The effect on recombination is similar in magnitude:
\begin{equation}
\frac{\Delta m_e}{m_e} \sim \bar{\psi}_{\rm rec} \sim 10^{-7}
\end{equation}

Again negligible.

\subsection{Could a Larger $\psi$ Help?}

For a significant effect on recombination, we need $\bar{\psi}_{\rm rec} \gtrsim 10^{-3}$. This requires either:
\begin{enumerate}
\item $D_0 \gg 0.017$ (ruled out by lab constraints)
\item A different $\psi$ potential that enhances $\bar{\psi}$ at $z \sim 1100$
\end{enumerate}

Option 2 is possible if $\psi$ has a potential with a feature at the recombination energy scale. But this would be an ad hoc addition with no DFD motivation.

\subsection{Agent 10 Conclusion}

\begin{center}
\fbox{\parbox{0.9\textwidth}{
\textbf{Finding 10.1:} DFD's $\psi$-dependent $\alpha$ variation at recombination is $\Delta\alpha/\alpha \sim 2\times10^{-7}$. The effect on the CMB is $\sim 10^{-6}$, completely negligible.\\[6pt]
\textbf{Finding 10.2:} Significant $\alpha$ variation ($\gtrsim 10^{-3}$) is needed to affect recombination, but L\'{o}pez-Honorez et al.\ (2024) showed this is incompatible with local constraints.\\[6pt]
\textbf{Finding 10.3:} $\psi$-dependent electron mass modifications are equally negligible ($\sim 10^{-7}$).\\[6pt]
\textbf{STATUS:} \RED{} --- Recombination modification cannot help. The $\psi$ amplitude at $z \sim 1100$ is far too small.
}}
\end{center}

\newpage
%=============================================================================
\section{Summary of Agents 6--10}
\label{sec:summary}
%=============================================================================

\begin{center}
\begin{tabular}{llcl}
\toprule
\textbf{Agent} & \textbf{Mechanism} & \textbf{Status} & \textbf{Key Finding} \\
\midrule
6 (Vector) & Add unit-timelike vector & \YELLOW & Works but costs 3 params + QG revision \\
7 (SymBoltz) & Numerical computation & \YELLOW & Essential next step; unlikely to save scalar-only DFD \\
8 (Nonlinear) & Mode coupling & \RED & Second-order $\sim 10^{-2}$; negligible \\
9 (Feedback) & Baryon-$\psi$ loop & \RED & Amplification $\sim 10^{-40}$; utterly negligible \\
10 (Recomb.) & $\alpha$ variation & \RED & $\Delta\alpha/\alpha \sim 10^{-7}$; negligible \\
\bottomrule
\end{tabular}
\end{center}

\textbf{Overall assessment:} Of the five resolution paths examined, only the vector extension (Agent 6) offers a viable path to fixing the third peak. All other mechanisms are negligible by many orders of magnitude. The SymBoltz.jl computation (Agent 7) is the critical next step to definitively quantify the deficit and validate the analytic estimates.

\end{document}
